\chapter{Limitations and Future Work}\label{chap:todo}

\section{Limitations}\label{sec:limitations}

Two boundaries define the scope within which these claims hold. The decision-layer results were obtained inside a synthetic-learner simulation rather than a real classroom, so they establish what the policy achieves against modelled learners, not yet against human ones. The perception-layer results were measured on the \ac{DAiSEE} benchmark rather than on live webcam streams, so they characterise the model under curated recording conditions rather than the variability of real deployment. The two layers, although joined by a single adapter interface, were never run together as one loop on real affect, so the end-to-end system is validated component by component rather than as a whole. These limitations do not undermine the results; they locate them. Each was also a deliberate methodological choice: a controlled benchmark isolates the perception model from confounded preprocessing, and a theory-grounded simulator isolates the policy from the ethical and logistical cost of experimenting on students before the architecture is ready to meet them. The directions in Section~\ref{sec:future_work} address each of these boundaries in turn.

\section{Future Work}\label{sec:future_work}

The first direction for future work is classroom deployment of the system. The current work is evaluated in a simulated environment, so testing it with real students is a natural next step. This is not only a practical step but a scientific one, because the policy was trained on simulator-derived affect with clean Gaussian noise and perfectly synchronous timing, whereas real facial emotion detection produces different noise and delays that the policy has never encountered. In real use, the system will face more noise in emotion detection, different student behaviours, and privacy constraints. To address this, the model can run on low-cost devices with on-device processing so that raw video does not need to be stored or shared. In addition, simple safety rules can be kept under teacher control so that educators can adjust system behaviour during use.

The second direction is to improve the perception part of the system using simple additional signals. Instead of relying only on facial expression, the system can also use low-cost modalities such as gaze direction, head movement, or speech tone (prosody). These signals are easy to collect and can improve robustness when facial emotion detection is uncertain. This matters most for the states the face hides: frustration and confusion reach an F1 of only $0.255$ and $0.224$, and these are exactly the cases that visual-only recognition misses. This extension does not change the main model, but only adds extra input signals to make predictions more stable.

The third direction is to make the system more flexible in classroom settings using a simple multi-agent setup. Instead of one learning agent only, a second agent can represent the teacher or classroom strategy. The two agents can work together: one focuses on student engagement, and the other controls teaching decisions. This makes the system easier to adapt to real classroom situations where teaching is not fixed and can change depending on student response.

Finally, future work can extend the system from short sessions to longer learning periods. This matters because the current policy reads a fixed-length state vector and cannot reason over long interaction histories, so slow-drifting states such as boredom are effectively invisible to it. In this case, the system would track student progress over multiple lessons instead of one session only. This would make the model closer to real education settings, where learning happens over time rather than in a single interaction.